\section{서론}\label{sec:intro}

스파이킹 신경망(Spiking Neural Network, SNN)의 효율성 서사는 하나의
대수적 사실에 기초한다. 뉴런 간에 오가는 것이 실수 활성값이 아니라
이진 스파이크이므로, 추론의 기본 연산은 곱셈-누산(MAC)이 아니라
\emph{발화한 곳의 가중치를 더하는 것}으로 충분하다는 사실이다
\cite{maass1997networks,roy2019towards}. 덧셈은 곱셈-누산보다 수 배
싸고\cite{horowitz2014computing}, 발화는 희소하므로 더할 일
자체가 드물다. 이 두 성질 --- 무곱셈성과 이벤트 희소성 --- 이 SNN이
약속하는 효율의 원천이며, Loihi나 TrueNorth 같은 뉴로모픽
프로세서는 이 약속을 실리콘 수준에서
구현한다\cite{davies2018loihi,merolla2014truenorth}.

문제는 이 약속이 \emph{배포 경로에서 소멸}한다는 데 있다. 이미
대규모로 보급된 범용 마이크로컨트롤러(MCU) 위에 SNN을 올리는 현실적
경로는 표준 TinyML 스택 --- TensorFlow Lite for
Microcontrollers(TFLM)\cite{david2021tflm}와
CMSIS-NN\cite{lai2018cmsisnn} --- 이고, 이 스택의 커널은 조밀(dense)
텐서 산술을 전제한다. 이진 스파이크는 int16 텐서로 실체화되어 dense
MAC 루프에 공급되고, 프로세서는 대부분이 0인 피연산자를 성실하게
읽어 와서 곱한다. 즉 모델은 이벤트를 말하는데 커널은 텐서를
계산한다. 저자의 선행 연구\cite{cha2026fusion}가 그 전형이다:
Conv-AvgPool-LIF 퓨전 연산자를 TFLM에 통합하여 표준 경로 대비 4.0배를
얻었지만, 커널 자체는 dense SMLALD 누산이었고 도달 처리량은 칩
이론 피크의 약 18\%에 머물렀다. 그 사이 실측된 은닉층 발화율은
4.1\% --- 지배 레이어 연산의 약 96\%가 0과의 곱셈이었다.

본 논문의 논제는 다음과 같다. \textbf{SNN의 이벤트 구동 이득은
하드웨어의 성질이 아니라 산술의 성질이므로, 하드웨어를 바꾸지 않고
커널 수준에서, 그것도 비트 단위로 동일하게 회수할 수 있다.} 입력이
$\{0, S\}$ 이진이면 내적은 다음과 같이 붕괴한다:
$\sum_k w_k x_k = S\cdot\sum_{x_k\neq0}w_k$. 이 항등식은 정수 산술의
결합$\cdot$교환법칙 위에서 성립하므로, 이를 구현한 커널은 근사
가속기가 아니라 \emph{의미를 보존하는 프로그램 변환}이다. 회수를
실제 성능으로 연결하기 위해 본 논문은 세 가지 설계 원리를 제시한다:
(P1) 활성 \emph{값}이 아니라 발화 \emph{위치}를 실체화하라 ---
im2col이 통째로 사라진다; (P2) 불규칙한 이벤트 축은 인덱스로,
규칙적인 채널 축은 연속 스트림으로 배정하라 --- 가중치 전치와
레지스터 블로킹이 그 귀결이다; (P3) 정수 에필로그를 불변으로 지켜라
--- 비트 동일성이 검증 가능한 계약이 된다.

본 논문의 두 번째 논제는 방법론적이다. \textbf{커널을 다시 쓰기
전에, 병목의 위치는 대상 워크로드 위의 절제(ablation) 실측으로
확정되어야 한다.} 사전 하드웨어 조사가 유망하게 지목한 두 가지 명령
페치 최적화 --- 플래시 가속기(ART) 활성화와 핫 커널의 ITCM-RAM 배치
--- 는 마이크로벤치마크에서는 각각 1.4배$\cdot$2.5배급 효과였으나,
실제 추론 워크로드에서는 $\le$0.12\%의 변화에 그친 실질 무효과였다.
우리는 그 원인(부팅 주소와 실행 주소의 불일치, 핫 루프와 캐시
용량의 관계)을 규명하고, 병목이 ``0을 곱하기 위한 데이터
이동''임을 사이클 수준 회계로 특정한 뒤에야 커널을 재설계했다.
이 진단 과정과 부정적 결과를 본문에 그대로 보고한다.

기여는 다음과 같이 요약된다.
\begin{itemize}
\item \textbf{개념화}: MCU 위 SNN 추론에서 이벤트 구동 이득이
  소멸하는 지점을 커널 경계로 특정하고, 그 회수를 세 가지 설계
  원리(P1--P3)와 의미 보존 계약으로 정식화한다(IV장).
\item \textbf{진단 방법론}: 조건을 레지스터로 증명하고 결과를 출력
  체크섬으로 불변식화한 사이클 정밀 절제 실측으로, 통념적 병목
  가설 2종을 반증하고 데이터 이동 병목을 확정한다(III장).
\item \textbf{구현과 실측}: 제안 커널은 TFLM 커스텀 연산자 프레임과
  CMSIS-NN 지원 함수 위에 그대로 얹히며, STM32F746 실측에서 종단
  $6.12$배($2{,}328.5\rightarrow380.6$ms/이미지), 발화율 4.1\%
  레이어 13.0배(동일 배치 dense 대비)를 여섯 빌드 변형의 출력 비트
  일치와 함께 달성한다(V장).
  런타임 오버헤드는 2.1\%에 머문다.
\item \textbf{일반화}: 실측에 적합시킨 1차 비용 모델로 희소-조밀
  손익분기 밀도를 약 89\%로 추정한다 --- SNN의 통상 동작 영역
  전체가 희소 커널의 우세 구간에 들어온다(VI장).
\end{itemize}

이후 구성은 다음과 같다. II장은 관련 연구 속에서 본 연구의 위치를
정한다. III장은 대상 시스템과 진단 방법론, 그리고 병목의 확정을
다룬다. IV장은 설계 원리와 그 구현$\cdot$비용 모델을, V장은 연구
질문 네 개에 대한 실측 답을 제시한다. VI장은 일반화와 한계를
논의하고 VII장에서 결론을 맺는다.
